\section{Further research questions}

\begin{enumerate}[leftmargin=*]
\item Can a canonical-$R$ encoding absorb zero-counter erasure into an
already necessary traversal and improve the current compact source minima,
while validating its complete public language? The normalization must be
applied after substitution, not merely when reading the original input.
\item Can the present normalization equivalence be formalized by extending
the existing Lean term and native-calculus development? Body extensionality
should be proved before normalization soundness; literal-source and
witness tests should be handled after unrestricted equivalence.
\item Which additional principle, beyond zero-equation induction, rules out
the two-point fibers of Section~\ref{zi:section}? Actual-equality reflection
suffices, but its cost must be counted. Does a narrowly stated
identity-iterator uniqueness principle already suffice for full equation
induction?
\item Can the new unary obstruction be strengthened to a useful normal
form or effective classification of the scalar self-seeded term functions?
The proof above supplies existence of a bound, not an algorithm extracting
an optimal bound from a term.
\item Can a pointfree proof representation reduce the substitution
machinery without introducing arbitrary extra pairing structure? External
tuple contexts avoid that interpretation-composite problem, but their
typing and conversion costs must be counted.
\end{enumerate}

\paragraph{Research status.}
The six-rule theorem, unary obstruction, and zero-induction countermodel are written mathematical
arguments. The executable checks test the stated implementations, scoping,
normalization, and selected derivations; they do not constitute a formal
verification of either universal metatheorem. The upstream formalized
scalar bridge is used with its stated scope. No global minimality or
literature-wide priority theorem is asserted.

